\section{Cross-fitted recalibrated PPI}\label{sec:method}
Proposition~\ref{prop:bound} suggests the obvious route to shorter intervals: estimate the calibration curve $g=\mathbb E[Y|Z]$ from the labeled data instead of fixing it to $z\mapsto\lambda z$. The estimator is an augmented-inverse-probability-weighted estimator with known labeling probability \citep{robins1994estimation}, equivalently a model-assisted/model-calibrated survey estimator with a nonparametric working model \citep{sarndal1992model,montanari2005nonparametric}; it is also the mean-estimation special case of cross-fitted DML \citep{chernozhukov2018dml}. We do not claim the estimator is new; the point of this paper is to measure how much it can gain over PPI/PPI++ and when.

\paragraph{Algorithm.} Split the labeled indices into $K=5$ random folds $L_1,\dots,L_K$. For each $k$ fit $\hat g_{-k}$ on $L\setminus L_k$ by a one-dimensional regression of $Y$ on $Z$. Put $\hat g_i=\hat g_{-k(i)}(Z_i)$ for $i\in L$ (out-of-fold) and $\bar g_j=\frac1K\sum_k\hat g_{-k}(Z_j)$ for unlabeled $j$. Then
\[
\hat\theta=\frac1n\sum_{i\in L}(Y_i-\hat g_i)+\frac1m\Big(\sum_{i\in L}\hat g_i+\sum_{j\notin L}\bar g_j\Big),\qquad
\hat V=\frac{\widehat{\mathrm{Var}}_L(Y-\hat g)}{n}+\frac{\widehat{\mathrm{Var}}_{\text{all}}(\hat g)}{m},
\]
and the interval is $\hat\theta\pm z_{1-\alpha/2}\sqrt{\hat V}$ (sample variances with $n{-}1$ and $m{-}1$ denominators). With $\hat g\equiv\lambda z$ this is PPI up to the use of the labeled predictions in the second mean; with $\hat g$ the OLS line it is the survey regression estimator, asymptotically equivalent to power-tuned PPI. Learners (all in \texttt{code/cppi.py}): \textbf{lin} (OLS line); \textbf{iso} (isotonic regression, increasing, \citealp{barlow1972isotonic}); \textbf{bin10} (10 equal-mass bins, i.e.\ post-stratification on the predictor output, in the spirit of stratified PPI \citealp{fisch2024stratified} but with random rather than allocated labels); \textbf{spl} (natural cubic regression spline with 5 knots at the 5\%--95\% quantiles of the training scores, \citealp{hastie2009elements}); \textbf{auto} (choose among lin, iso, spl by inner 4-fold CV mean squared error, nested inside each outer fold, so that selection never sees the points it is evaluated on). No tuning on the evaluation data was done; $K$, knot counts and bin counts were fixed before the experiments.

\begin{proposition}[Validity]\label{prop:valid}
Suppose $\mathbb E Y^2<\infty$ and there is a function $g^*$ with $\mathbb E g^*(Z)^2<\infty$ such that $\mathbb E[(\hat g_{-k}(Z)-g^*(Z))^2\mid\text{training fold}]\to_p0$ for each $k$, where $Z$ is an independent copy. Then $(\hat\theta-\theta)/\sqrt{\hat V}\Rightarrow N(0,1)$, so the interval has asymptotic coverage $1-\alpha$. Moreover $\hat\theta$ is exactly conditionally unbiased given the training folds for \emph{any} $\hat g$. If $g^*=g$ the interval attains $V_{\mathrm{eff}}$.
\end{proposition}
The limit $g^*$ need not be the truth: a misspecified learner (e.g.\ a monotone fit to a non-monotone $g$) yields a valid but longer interval, with variance $\mathbb E(Y-g^*)^2/n+\mathrm{Var}\,g^*(Z)/m$. This is the same robustness-to-the-predictor that makes PPI valid for any $f$; the proof (Appendix~\ref{app:proof}) uses only conditional Chebyshev bounds and no convergence \emph{rate} for $\hat g$. The guarantee is asymptotic; Sec.~\ref{sec:exp} measures finite-sample coverage down to $n=25$. Data-driven learner selection (\textbf{auto}) is covered by the proposition when the selected learner's limit is stable; we do not prove this and rely on simulation.
